\chapter{Background and Literature Review}
\label{ch:background}

\section{Retrieval as part of an answering system}
A language model can produce a fluent answer without exposing where its
information came from. Retrieval offers a different interface: relevant
external text is selected and made available to the answering model.
Lewis et al.\ formalised a RAG architecture combining a retriever with a
sequence-to-sequence generator, bringing non-parametric information into
the prediction process \cite{lewis2020rag}. In this thesis, that broader
architecture motivates the task, but only the evidence-selection stage
is evaluated.

Retrieval itself has several decisions. The system needs a searchable
representation, a relevance function, a candidate set, and a rule for
constructing the final context. Granularity affects all four. Changing
the boundaries of a passage changes its embedding, the number of indexed
items, the competition between candidates, and the amount of text returned
for a fixed number of results. It is therefore more than a formatting
choice made after retrieval.

It is helpful to distinguish document retrieval from passage retrieval.
In an open-domain setting, finding the relevant document is part of the
task. In this study, the associated QASPER paper is known and retrieval
is filtered to that paper. The remaining difficulty is selecting evidence
inside it. A method that succeeds under this restriction may still fail
when unrelated documents are introduced, because the score distributions
and the nature of hard negatives change.

\section{Dense representations and similarity search}
Dense retrieval maps questions and passages into a vector space. A
similarity function then ranks passages by their relation to the question.
Dense Passage Retrieval demonstrated the effectiveness of learned
question and passage encoders for open-domain question answering
\cite{karpukhin2020dpr}. The present project uses a fixed embedding
service rather than retraining such a retriever. Consequently, any
improvement must come from selecting or organising the existing
representations, not from learning a better embedding space.

For nonzero vectors $\mathbf{e}_q$ and $\mathbf{e}_c$, cosine similarity is
\begin{equation}
s(q,c)=\frac{\mathbf{e}_q^\top\mathbf{e}_c}
{\|\mathbf{e}_q\|_2\|\mathbf{e}_c\|_2}.
\label{eq:cosine}
\end{equation}
The repository uses \texttt{text-embedding-3-small} with 1,536-dimensional
vectors \cite{openai2024embeddings}. It keeps the model and dimension
fixed across the reported retrieval comparisons. This controls an
important source of variation, although it also limits the conclusions
to one embedding representation.

Qdrant stores vectors together with payload fields used for filtering,
including paper identity and granularity \cite{qdrantcollections}.
Its vector index is an infrastructure component, not the document
hierarchy proposed later in the thesis. HNSW organises vectors for
approximate nearest-neighbour search \cite{malkov2016hnsw}; the
similarity tree in this work organises aligned text spans within a
paper. The two structures solve different problems despite both using
hierarchical organisation.

A high cosine score should not be read as a calibrated probability that
a passage contains sufficient evidence. In particular, scores produced
for different passage lengths may have different distributions. The
mixed-size experiments directly test what happens when those scores
are pooled without learning a cross-granularity calibration. Their
outcome is therefore informative about this pipeline, but it is not a
general theorem about dense retrieval.

\section{What retrieval granularity changes}
\subsection{Precision, coverage, and context}
A small chunk can concentrate on words closely related to a question.
It can also cut away a definition, a qualifier, or the subject of a
pronoun. A large chunk can retain those connections while bringing in
text that does not belong to the evidence. Neither choice is uniformly
better. The right balance depends on the scoring objective as well as
the document.

In a top-$K$ evaluation, length and count interact. With five full chunks,
10-token retrieval returns approximately 50 source tokens, whereas
160-token retrieval returns approximately 800. The two methods have
the same number of results but not the same text budget. A rise in recall
at the larger size may simply reflect this increase in coverage.
Likewise, a decline in precision may reflect the additional denominator
in the lexical metric. A fair interpretation must retain both facts.

Evidence annotations introduce a further distinction. The number of
tokens highlighted by an annotator describes the reference material,
not necessarily the best size of an independently ranked chunk. If
evidence consists of several short statements, retrieving several
compact chunks may be preferable to retrieving several long chunks
whose length resembles the total annotation. This observation motivates
the later move from evidence-length classification to retrieval regret.

\subsection{Fixed windows and semantic units}
Fixed token windows are simple and reproducible, but their boundaries
need not match sentences or propositions. Dense X Retrieval studies
retrieval units at several levels and argues for propositions as atomic,
self-contained pieces of information \cite{chen2024densex}. Its units
are constructed semantically. A 10-token window in this project is not
equivalent to a proposition, and the thesis does not claim to reproduce
that method.

The advantage of the fixed windows here is experimental alignment. The
chosen sizes double at each level, so a coarse window contains smaller
windows from the same token sequence. This makes parent--child score
comparisons possible without an additional sentence segmentation or
semantic clustering procedure. The disadvantage is equally concrete:
the resulting hierarchy preserves positional containment, not
necessarily a meaningful discourse structure.

An alternative way to retain context is to change when embeddings are
constructed. Late Chunking encodes a longer text before pooling token
representations into chunk embeddings, allowing the chunk representation
to retain surrounding context \cite{gunther2024late}. The repository
instead embeds each stored chunk as its own text. The current results
therefore do not answer whether adaptive routing remains useful when
small chunks already have contextualised representations.

\section{Hierarchical retrieval}
Hierarchical approaches offer access to information at more than one
scale. RAPTOR recursively groups and summarises text to build a
tree-organised retrieval structure \cite{sarthi2024raptor}. Its upper
levels contain generated summaries. The hierarchy in this thesis is
different: a 160-token node is a literal span of the paper, not an
abstraction generated from its descendants. This avoids an extra
summarisation model and preserves direct source offsets.

The distinction matters for both capability and evaluation. A generated
summary can express information not present in any single child and can
combine evidence across separated parts of a document. A contiguous
token tree cannot do that. On the other hand, its parent and child
spans have a transparent containment relation, making it possible to
measure local overlap with annotated evidence and to inspect score
differences without assessing summary faithfulness.

The similarity-tree hypothesis is that these score differences reveal
something about evidence concentration. A strong parent and consistently
strong children may suggest a broadly relevant region. A single strong
child among weak siblings may suggest a narrow match. This is an
interpretation of possible signals, not an assumption that either
pattern always identifies gold evidence. The experiments test whether
a classifier can use these patterns more effectively than it uses
level-wise score summaries alone.

\section{Adaptive retrieval in RAG}
Adaptation can occur at several points in a retrieval pipeline.
Adaptive-RAG uses question complexity to select among answering
strategies with different retrieval demands, including different
amounts of retrieval activity \cite{jeong2024adaptive}. Its decision is
about the retrieval strategy, rather than choosing among the five
fixed token lengths used here. The shared idea is that a question can
inform how resources are allocated; the action spaces and supervision
are different.

Self-RAG integrates decisions to retrieve with generation and
self-reflection through special reflection tokens
\cite{asai2023selfrag}. It addresses the interaction between retrieval,
generation, and critique. This thesis keeps generation outside the
evaluation and makes a narrower decision before evidence is returned.
A successful router here would therefore be a possible retrieval
component, not a replacement for a system that verifies its generated
claims.

Mix-of-Granularity (MoG) is a particularly close reference point for
the present study. It uses a learned router to select retrieval
granularity from the input query instead of applying one fixed
granularity to every query, and trains that router with soft-label
supervision \cite{zhong-etal-2025-mix}. Its graph extension,
MoG-Graph (MoGG), reorganises documents as graphs so that relevant
pieces that are non-adjacent in the original text can be retrieved
through graph neighbourhoods. The present thesis does not reproduce
MoG or MoGG. It instead studies evidence retrieval within a known
QASPER source paper, using aligned fixed GPT-2 token spans at sizes
$\{10,20,40,80,160\}$. Its primary outcome is joined evidence-token
retrieval F1 rather than generated-answer accuracy, and it compares
evidence-length classification, question-only Qwen routing,
similarity-tree features, and direct retrieval-regret supervision.

The amount of supplied context also cannot be separated entirely from
how a generator uses it. Lost in the Middle reports sensitivity to the
position of relevant information in long contexts
\cite{liu2023lost}. This motivates care in assembling evidence, but
does not establish that a higher lexical retrieval F1 will produce
better answers in the present setting. That connection requires a
downstream experiment with a specified generator and context ordering.

\begin{table}[tbp]
\centering\small
\caption{Position of the thesis relative to neighbouring retrieval ideas.
The rows describe different decisions, not a ranking of methods.}
\label{tab:literature}
\begin{tabularx}{\textwidth}{@{}p{3.1cm}XX@{}}
\toprule
Approach & Main decision or representation & Difference from this thesis\\
\midrule
Dense X Retrieval & Semantic retrieval units, including propositions & Fixed token spans are not semantic propositions\\
RAPTOR & Hierarchy with recursively generated summaries & Trees here contain literal, aligned source spans\\
Late Chunking & Contextualise before forming chunk embeddings & Stored chunks here are embedded independently\\
Adaptive-RAG & Adapt the answering/retrieval strategy & The primary action here is chunk size\\
Self-RAG & Retrieve, generate, and critique jointly & Only evidence retrieval is evaluated here\\
Mix-of-Granularity (MoG) & Query-dependent routing across retrieval granularities & Medical QA; here, compare routing targets, including retrieval utility, over aligned token spans within a known QASPER paper\\
This work & Select global or local token granularity & Compare decisions under fixed embeddings and explicit evidence scoring\\
\bottomrule
\end{tabularx}
\end{table}

\section{Learning a router}
\subsection{Classification and its target}
A router can be expressed as a classifier over five sizes. Given features
$\mathbf{x}$, it predicts probabilities $p_\theta(g\mid\mathbf{x})$ and
chooses the largest. Standard cross-entropy rewards the labelled size
and penalises every other size through the same categorical objective.
Its suitability depends on how that label was obtained.

If the target is evidence length, a correct classification means that
the model has recovered a length category. It does not mean that the
associated top-five retrieval set is optimal. If two sizes have almost
identical retrieval utility, a hard label also hides that similarity.
Conversely, two equally wrong classifications can have very different
retrieval costs. These are objective-design issues, not simply limits
of classifier capacity.

Class imbalance creates an additional difficulty. Class-balanced loss
uses the effective number of samples to reduce the dominance of common
classes \cite{cui2019classbalanced}. This provides the weighting
principle tested in Phase 2B. It cannot, by itself, make a proxy target
more faithful to the downstream retrieval objective.

\subsection{Nonlinear models and combined features}
Boosted decision trees provide a way to learn nonlinear interactions
between score statistics without fine-tuning the embedding model.
XGBoost supplies a regularised gradient-boosting implementation
\cite{chen2016xgboost}. Here its role is modest: learn whether particular
combinations of level-wise and hierarchical features support a size
decision. Feature importance can suggest which measurements the fitted
model used, but it does not establish a causal explanation of relevance.

Combining a language model and a tree model raises a separate training
question. If a second-stage learner sees base-model predictions on
examples that trained the base model, its inputs can be unrealistically
confident. Out-of-fold predictions are a practical way to reduce this
mismatch. Grouping by paper is important because questions about the
same paper share content and should not be treated as independent
documents when constructing the folds.

\section{Evaluation as part of the method}
Evaluation is not a final reporting step detached from model design.
The target, decision unit, candidate set, and retrieval budget determine
what a score means. A question-level classifier, a local-tree classifier,
and a retrieval system can all report an F1 value while measuring
different events.

Uncertainty also depends on the unit of observation. Questions from the
same paper share a source, so resampling individual questions can ignore
an important dependence. The later repository comparisons resample
papers and preserve all of their questions. General guidance on
significance testing in NLP emphasises choosing procedures that match
the task and comparison \cite{dror2018testing}. The exact paired
paper-cluster bootstrap used here is specified in
Chapter~\ref{ch:results}.

The gap addressed by this thesis is consequently specific. It is not
that granularity or adaptive retrieval has never been studied. It is
the empirical question of which signals and learning targets support
a useful granularity choice in this fixed, evidence-annotated setting,
and whether their added complexity beats a strong fixed-size baseline.
